\begin{proof}
\begin{enumerate}[label=\arabic*.,leftmargin=*]
\item For $n=1$ first blow up $(u,q)$ and then the pullback
of $(f,g)$. The charts of the resulting modification are
\[
\begin{array}{c|c|c}
 \text{coordinates}&\text{map to }X_1&r\\\hline
 (q,U,v,f)&u=qU,\ g=fUv&f(Uv-1)\\
 (u,Q,G,f)&q=uQ,\ v=GQ,\ g=fG&f(G-1)\\
 (u,v,F,g)&q=uvF,\ f=gF&g(1-F).
\end{array}
\]
To derive them, the $q$ chart of the first blow-up
has $u=qU$, hence $g=fUv$; the ideal $(f,g)$ is
already principal there. Its other chart has
$q=uQ$ and equation $fv=gQ$. Blowing up $(f,g)$
gives the last two charts. These are the 
blow-up charts, with the common factors cancelled
on the closure of the dense torus; no additional
component of the naive total transform is included.
They are smooth. In the first chart $Uv-1$ is a
smooth coordinate along its zero set, and in the
others $G-1$ or $1-F$ is a coordinate. Thus the
displayed products prove the reduced NC assertion
for $r$.

\item For arbitrary $n$, pull back this proper modification
by $q=s^n$. The first chart remains smooth. The
second is $uQ=s^n$ times the two coordinates $G,f$.
Resolve its $A_{n-1}$ factor by the unit-interval fan.
The third chart is
\[
 uvF=s^n
\]
times the independent coordinate $g$. Its cone lattice
may be written as $N=\mathbb Z^3$ with coordinates
$(a,b,d)$, where the valuations of $u,v,F,s$ are
$a,b,nd-a-b,d$. The cone is the cone over
$n\Delta_2$. Insert all rays $(i,j,1)$ with
$i,j\ge0$ and $i+j\le n$, and triangulate by the
unit triangles
\[
\begin{aligned}
 &(i,j),(i+1,j),(i,j+1) &&(i,j\ge0,\ i+j\le n-1),\\
 &(i+1,j),(i,j+1),(i+1,j+1)
                       &&(i,j\ge0,\ i+j\le n-2).
\end{aligned}
\]
Every determinant is one in absolute value; the
triangles partition the large triangle and give a
proper smooth toric subdivision. Their number is
$n^2$. All new rays have $s$-height one.

\item The open set $F\ne0$ is the toric open corresponding
to $nd-a-b=0$. Its fan is the unit subdivision of
the edge $i+j=n$, hence the minimal $A_{n-1}$
resolution times $\mathbb C^*$. This is precisely
the second chart's resolution under $Q=vF$.
Where that chart meets the first, the $A$-type
resolution is an isomorphism. These facts glue
the local proper resolutions to a proper
modification of the whole local $X_n$.

\item The function $F$ is a toric monomial. Above $F=1$
it is submersive: in a regular toric chart a
nonzero monomial has nonzero derivative in a
coordinate with positive exponent. Equivalently,
on $F\ne0$ it is the displayed torus projection.
Thus $r=g(1-F)$ has two smooth transverse
branches in the third chart after resolution.
The second chart has $r=f(G-1)$ independently
of the resolved surface coordinates, and the
first was already checked. This proves smoothness
and reduced normal crossings over $r$ for all $n$.

\item For $r\ne0$ in the third chart,
$g=r/(1-F)$, so the local fibre is the resolved
threefold $uvF=s^n$ with $F\ne1$.
For $n=1$, its $q$ divisor is $uvF=0$, a literal
three-branch normal-crossings chart. For $n=12$
the standard triangulation gives 144 unimodular
charts above this one triple point. This is the
local cyclic-cover count at each cusp triple point. The other
potential singularity, at $f=-r$, is included
in this construction: the first blow-up small-resolves it
for $n=1$, and the second is an isomorphism there
because $f$ is nonzero.

\item The equation arises from the modified neighbourhood of \ref{push:src:B.4.21}. There $h=uv-q$ in the
node chart before modification, and $H=h/q$ is holomorphic
in the modified collar. The unit inverse-normal plumbing
$Hf=r$ is exactly
\[
 f(uv-q)=rq,
 \qquad fuv=(f+r)q.
\]
Putting $g=f+r$ gives $X_n$, including its factor $q$.
The modified collar, where $H$ is small,
identifies with the third chart near $F=1$ by
\[
 F=(1+H)^{-1},\qquad g=f(1+H),\qquad r=fH.
\]
There $uvF=s^n$ is exactly the resolved equation
$uv=s^n(1+H)$ from \ref{push:src:B.4.21}. The charts near $F=0$
are additional local completion charts.


\item The first collar action is $H\mapsto aH$, $f\mapsto a^{-1}f$,
which fixes $r=fH$. The function $G(q,H)$ of \ref{push:src:B.4.21}
conjugates this action on its stated small collar. For fixed $q$,
the product converges locally uniformly for $H\in\C$, because
$\sum_j|a(q)^jH|$ does so on compact sets. At $H=-a^{-j}$ it
vanishes, and the coordinate change $Z=xG(q,H)$ has zero derivative
in $x$. For $b\in\C^*$, on the open set where $r\ne0$, $f\ne0$ and
$W=1+r/f\notin\{1,b^{-1}\}$, the map $f=r/(W-1)$ conjugates
multiplication $W\mapsto bW$ to the rational map
\[
 f\longmapsto\frac{rf}{(b-1)f+br}.
\]
Indeed substituting $b(1+r/f)$ into $r/(W-1)$ gives the displayed
formula. This is the exact coordinate comparison for the test variable $W$.
\end{enumerate}
\end{proof}